\subsection{Overview}

Edge-of-field treatment structures---denitrifying woodchip bioreactors and
their companions (constructed wetlands, saturated riparian buffers, and
phosphorus-removal structures)---have accumulated abundant performance
demonstrations yet remain defined more by open design questions than by settled
rules. Bioreactors remove nitrate by microbially mediated denitrification, but
their performance is bounded by low water temperature, variable hydraulic
retention time (HRT), and the biogeochemical evolution of the woodchip medium as
it ages. The frontier clusters into interlocking lines: how to size, scale, and
standardise design; how to resolve the internal hydraulics and tracer methods on
which design rests; whether the microbial engine can be deliberately enhanced;
how to select, amend, and manage reactive and phosphorus-sorbing media; how to
characterise multi-decadal aging, clogging, and construction effects; how
pumped and coupled configurations extend the design space; and how companion
practices complement or substitute for bioreactors. A recurring meta-message is
the need for standardised design procedures, broader performance-metric
reporting, and long-term data across climates, landscapes, and maturation stages
so that watershed-scale nutrient-reduction goals can be met
\parencite{Christianson2021,Maxwell2022,Bell2015}. The evidence is drawn
overwhelmingly from site-scale field and laboratory studies concentrated in the
U.S.\ Midwest; the by-products of reductive treatment (greenhouse gases,
dissolved organic carbon export, methylmercury) are treated here only as design
constraints, their downstream accounting deferred to the consequences synthesis.

\subsection{Sizing, scaling, and standardised design}

A central frontier is reconciling site-level effectiveness with watershed-scale
load removal. \textcite{Christianson2026} argue for better balancing bioreactors
that are as effective as practical at the site while still removing as much load
as possible; they flag the untested performance of \emph{undersized} bioreactors
serving larger drainage areas, propose greater design flexibility to scale the
practice, and caution that additional empirical data are needed to strengthen the
loading-density relationship and that performance data from full-size bioreactors
beyond Illinois are required before prediction models are used for design. For
hydrologically flashy sites, \textcite{Dougherty2020b} propose sizing designs
upward and repeatedly call for simpler, standardised ``plug-and-play'' design
procedures; standardisation is echoed at the policy level by
\textcite{Christianson2021}, who recommend refining the Conservation Practice
Standard and streamlining the design process. The adequacy of residence-time
($T$)-centric design is itself questioned: \textcite{deOliveira2023} advocate
field-scale assessment of baffles in wide bioreactors to raise hydraulic loading
while preserving efficiency, and more tracer testing of non-standard media
coupled to nitrate monitoring and modelling. Foundational validation remains
incomplete---further field-scale validation of the design procedure against beds
of differing design \parencite{Christianson2013}, validation of hydraulic models
at additional denitrification-bed sites with conservative isotopic tracers
\parencite{Ghane2015}, region-specific parameterisation of flow, seasonality, and
microbial limits for new settings such as Lithuania \parencite{Povilaitis2020},
and sorption experiments underpinning tracer-based design \parencite{Ghane2019}.

\subsection{Hydraulics, retention time, and tracer methods}

HRT is repeatedly identified as the master variable linking hydraulics to
removal, yet its treatment in design guidance is imperfect. \textcite{Moorman2015}
note that research on the HRT--nitrate-removal relationship would sharpen sizing
criteria, and \textcite{Povilaitis2018} call for proper optimisation of HRT
effects through both laboratory- and field-scale study across flow conditions. In
a global synthesis, \textcite{Christianson2025c} recommend methodological
advances: assessing varying HRT by changing flow rate versus sampling along the
reactor length; greater transparency in estimating woodchip porosity; more
consistent reporting of dissolved oxygen, pH, and dissolved organic carbon; and
machine-learning refinement of HRT--removal-efficiency estimates as data
accumulate. The internal hydraulics of woodchip media are poorly resolved:
\textcite{Ghane2016} argue that the non-Darcy flow characteristics of woodchips
should be exploited to revise NRCS design guidelines. Tracer methodology is a
specific bottleneck. \textcite{Ghane2019b} report anomalously low bromide
recovery and recommend testing for organobromine in woodchips, microscopic and
excavation-based examination of physical bromide entrapment, and either
accounting for retardation or adopting non-sorbing conservative tracers such as
isotopes; \textcite{Ghane2018} likewise recommend tracer tests to compare the
hydraulic performance of retrofitted beds. Reliable tracer methods are thus a
prerequisite to benchmarking hydraulic performance at all.

\subsection{Enhancing the microbial engine}

The defining enhancement frontier is whether denitrifiers can be inoculated into
woodchip bioreactors to raise removal where cold temperatures limit biological
activity. \textcite{Jang2019} propose that bioaugmentation might increase
low-temperature removal but note that genome completion is needed to identify the
nitric-oxide reductase of a candidate \emph{Cellulomonas} inoculant.
\textcite{Anderson2020} frame the first-order uncertainties: whether inoculated
strains enhance field removal, how long they survive in bioreactors, and whether
carbon addition boosts both inoculated and indigenous denitrifiers. Survival and
persistence recur globally, with \textcite{Wang2025b} calling for field-scale
methods that minimise inoculum loss and for exploration of emerging targets and
consortia---nitrate-dependent anaerobic methane oxidation (n-DAMO), biostimulation
to favour inoculated bacteria, and algal--bacterial granular sludge. Terminal
denitrification steps in fungal strains, including confirmation of nitrous-oxide
reductase and reduction of nitrate to dinitrogen under bioreactor conditions,
remain to be verified \parencite{Aldossari2021}, and microbial amendments require
replicated pilot-scale testing before upscaling
\parencite{Jeglot2022,Ma2021}. Against this enthusiasm,
\textcite{Hathaway2015} offer a testable counter-hypothesis---that optimisation is
more likely from refinements to design and management parameters that shape
environmental conditions than from direct microbial manipulation---a useful null
hypothesis for field trials.

Rational manipulation presupposes mechanistic understanding of the resident
community, which is incomplete. \textcite{Hartfiel2022} argue that bioreactor
microbiome study should draw on soil and wetland ecology; that controls on
cellulose- and lignin-metabolism genes warrant attention; that dissimilatory
nitrate reduction to ammonium and anammox have been observed but understudied;
and that varying substrates to select for differing communities is a promising
route to more consistent performance and reduced by-product formation---the
mechanism by which HRT shapes the community, and thereby by-product production,
being a specific gap. Complementary needs include managing bed pH for acidic
drainage and conceptualising regional controls on in-bed nitrogen removal
\parencite{Bock2018}, routine monitoring of nitrite to avoid its accumulation in
effluent \parencite{Kouanda2021}, and characterising the selective pressure of
oxygen versus nitrate load on community assembly \parencite{Schaefer2021}. The
emissions accounting of these by-products is treated in the consequences
synthesis; here they define a design and monitoring constraint.

\subsection{Reactive media, carbon sources, and structural materials}

Beyond conventional woodchips, alternative fill and carbon-management strategies
are incompletely characterised. \textcite{Law2023} encourage field studies of
corncobs capturing spatial and temporal variability, the nitrogen-removal-rate
range, and whether a lower maximum HRT is needed, while
\textcite{Porter2015} call for study of organic-carbon availability at different
depths and moisture contents, so that substrate choice or internal water levels
could be tuned. Because woodchip carbon fuels denitrification, its release
kinetics and pre-treatment are active frontiers: the longevity of
composting effects and the biodegradability and release rates of woodchip DOC
remain to be quantified \parencite{Kouanda2021}, and startup carbon losses raise
questions about recycling the initial flush or pre-washing chips before
installation \parencite{Abusallout2017}. Synthetic-carbon dosing is unresolved---
the source, frequency, and amount needed to optimise performance under dynamic
inflow while avoiding excess DOC and biofilm clogging \parencite{Rivas2020}, and
the hybrid application of synthetic carbon in low-temperature-tolerant reactors
\parencite{Zhang2021b}.

A distinct materials frontier concerns drain envelopes and structural fill.
Rice-husk envelopes require field testing of combination ratios and long-term
assessment given husk degradation and acidification
\parencite{Kaboosi2012,Navabian2017}; aggregate sizing needs in-situ testing in
clay soils \parencite{Byrne2023}; and synthetic geotextile envelopes must be
evaluated individually because of varying physical properties, with
filter-sock-plus-aggregate combinations not yet recommendable
\parencite{Byrne2023b}. Structural concrete alternatives---fly-ash- and
coal-gangue-based concrete---remain untested for durability under high
groundwater \parencite{Wang2016c}.

\subsection{Phosphorus-removal media and filter design}

Phosphorus-sorption materials (PSMs) constitute one of the most heavily populated
gap areas, spanning selection, longevity, regeneration, and hydraulic
integration. \textcite{Li2020} argue that adsorbent choice must weigh phosphorus
concentration, interfering ions, temperature, flow rate, and sediment load;
prioritise selectivity at neutral pH; be testable for reuse and recovery; and be
economical at field scale. \textcite{Mendes2022} add that the apparently
synergistic retention mechanism of filter media at high flow rates is unknown and
that regeneration of spent media should be pursued. Iron- and steel-based media
dominate specific proposals: optimising an iron-oxyhydroxide regeneration process
and evaluating low-cost goethite such as bog iron ore \parencite{Allred2017};
determining minimum contact time for zero-valent iron via higher-flow column and
isotherm tests \parencite{Allred2014}; regenerating steel by-products by alkaline
washing \parencite{Hua2016}; connecting a steel-slag filter to a bioreactor only
after the initial DOC flush and mixing fine steel to preserve conductivity
\parencite{Sellner2019b}; optimising removal, retention time, slag mass, and cost
\parencite{Wang2020d}; controlling sediment influx and improving flow
distribution \parencite{Shedekar2020}; and recognising that ligand-exchange PSMs
suit tile drainage only where structures allow high-frequency media replacement
(every 4--6 months) \parencite{Penn2020}. Fly-ash and pellet media raise
saturation-timing questions---saturation time must be quantified before field use
so pellets are replaced before shifting from phosphorus sink to source
\parencite{Li2017,Li2018b}---and filter blends remain unoptimised with respect to
media configuration, contact time, and high-flow contact
\parencite{Agrawal2011,Agrawal2011b,Kumari2025,Kankarej2023}. Integrated
woodchip-plus-PSM designs need multi-scale, multiyear evaluation of removal
effectiveness, media placement, lifespan, and regeneration \parencite{Husk2018}.

Novel phosphorus dynamics are emerging. \textcite{Zhang2024b} call for
experiments on enhanced orthophosphate uptake by woodchips undergoing
anoxic--oxic cycling and identification of polyphosphate-accumulating organisms,
while redox-driven phosphorus release remains a monitoring gap requiring start-up
flush and aging assessment with high-frequency dissolved-reactive-phosphorus
(DRP) monitoring and drawdown plates \parencite{Bailon2025}, multi-season
confirmation of DRP removal \parencite{Rivas2020}, more reactive passive-sampling
resins \parencite{Li2020b}, and characterisation of dissolved organic phosphorus
lability and bioavailability \parencite{Liu2024b}.

\subsection{Aging, longevity, clogging, subsidence, and construction}

The multi-decadal service life and physical stability of treatment structures are
among the least-studied dimensions. \textcite{Ghane2018} recommend longer study
of woodchip particle-size reduction, investigation of whether woodchip carbon
becomes recalcitrant to the point of inefficiency before particle-size reduction
throttles flow, and studies beyond ten years. \textcite{Christianson2020}
similarly call for long-term documentation of biofilm formation, siltation, wood
decay, and the effects of a decreasing carbon-to-nitrogen ratio and increasing
lignocellulose index, to inform woodchip-recharge thresholds; long-term
($>$8~yr) field programmes were underway to establish efficacy
\parencite{Greenan2009}, and spatiotemporal removal should be evaluated across
normal, wet, and drought years alongside woodchip C:N measurements
\parencite{Hartfiel2023b}. Estimating replacement interval and cost depends on
performance, lifespan, and residue cost, with pre-measured feedstock volumes
proposed to control construction-method effects \parencite{Kohn2023}. Chemical
clogging is mechanistically unresolved and needs simulation models for prevention
\parencite{Guo2020,Guo2021}. Subsidence monitoring is a new methodological
frontier: \textcite{Christianson2025b} propose applying a LiDAR method at
additional sites---especially immediately after construction---to build
subsidence time series and assess targeted inlet recharge. Sediment management at
construction is itself a design frontier: \textcite{Feyereisen2023} plan to weigh
the opposing effects of bypass flow and nitrate removal in a sediment-exclusion
system and advise pre-construction investigation of sedimentation sources.
Network-scale diagnostics also need extension---higher-pressure air injection and
time-domain electromagnetic resistivity monitoring for large drainage networks
\parencite{Henine2021}.

\subsection{Pumped, energy-assisted, and coupled configurations}

Pumping to raise hydraulic loading is an emerging but incompletely validated
avenue. \textcite{DíazGarcía2025} call for assessing different pumping
configurations, note that continuous 24-h operation or a longer pumping window
could increase treated flow, and caution that pumping may suit only the most
appropriate applications. \textcite{Hartfiel2024} advocate solar-powered pumped
bioreactors to reduce operating and unit cost, and \textcite{Hartfiel2023}
recommend sizing pumping to a bioreactor's actual rather than design capacity,
timing pumping to source-water nitrate concentration, temperature, and carbon
content, and extending applications to nitrate-laden surface waters, irrigation
waters, ditches, cisterns, and groundwater. Management experiments---partial
inlet-recharge strategies and dual-chamber operation engaging a second chamber
only under high flow---are proposed \parencite{Hartfiel2023b}, and pumped
constructed-wetland-plus-woodchip-bioreactor systems require analysis of
particulate versus dissolved nutrient fractions during intense discharges and
longer monitoring of the pumping effect on phosphorus retention
\parencite{Pugliese2020}. Novel source waters and enhanced configurations broaden
the space: denitrification responses to spring flow and concentration, with
payment-based incentives for spring bioreactors \parencite{Easton2019};
low-temperature bioelectrochemical reactors with horizontal flow, baffled
oxidising/reducing zones, and plate cathodes \parencite{Law2018}; and coupled
phosphorus-removal structures sited at critical source areas, including a proposed
redox-coupled Fe(II)--nitrate reaction for nitrate removal by activated aluminium
\parencite{Penn2024}.

\subsection{Wetlands, saturated buffers, and companion structures}

Frontier questions extend to practices deployed alongside or instead of
bioreactors. For saturated riparian buffers (SRBs), a key challenge is predicting
nitrate removal well enough to design more effective buffers, motivating
unsteady-flow and numerical-transport methods and optimal-width approximations
\parencite{McEachran2020}; replicating SRB results across landscapes, with
biomass harvesting to sustain viability \parencite{Jaynes2014}; more
sophisticated removal and hydrology components coupled to economic objective
functions, and determination of maximum allowable buffer length before optimising
width \parencite{Abdalaal2023}; confirmation that removal increases as vegetation
matures, with multipipe configurations and additional monitored sites
\parencite{Jaynes2019,Jaynes2019b,Singh2024b}; and paired-field determination of
actual performance and the cost-effectiveness of the buffer-flow and
controlled-drainage components \parencite{Ghane2025}.

Constructed and treatment wetlands present persistent performance-limit
questions. Hourly-time-step dynamic modelling is needed to approximate nitrate
removal in tile-fed wetlands with means to buffer and gradually release flows
\parencite{Tanner2011}, and long-term study is required of sustainable
nitrogen-retention efficiency across climates, whether organic-carbon
availability limits denitrification, the value of biomass cutting or slow-release
carbon supplements, and flow-regime effects analysed through residence-time
distributions \parencite{Tanner2005,Tanner2003,Tanner2005b}. A broad agenda for
constructed treatment systems spans driving factors, flow--treatment
relationships, design-parameter effects, long-term field tests, and
winter-efficiency enhancement \parencite{Mendes2021}, alongside maturation of
constructed wetlands for phosphorus retention \parencite{Mendes2020}, siting and
sizing by wetland-to-catchment ratio and seasonal discharge
\parencite{Steidl2019,Tournebize2017}, and subtropical adaptation of
land-treatment (FILTER) systems \parencite{Gardner2001}.

Denitrification-bioreactor and combined-installation refinement remains open:
reactor size, packing materials, elimination of preferential flow, and pH
maintenance \parencite{Aide2016}; separating the contributions of controlled
drainage with and without a bioreactor and evaluating clogging potential
\parencite{Husk2017}; and verification of design retention times and hydraulic
conductivity \parencite{Feyereisen2015}. Controlled and flow-regulating
structures raise further questions---improving the connection of regulator
elements to limit effective flow sizes \parencite{Popek2021}, pipe-riser and
event-scale trapping effects in water-and-sediment control basins
\parencite{Gupta2019}, and soil-surface clogging, storage geometry, and
discharge-capacity effects in drainage-inlet flow control \parencite{Nédélec2010}.
Finally, subsurface attenuation adjacent to drainage needs broader
characterisation: hyporheic-zone processes and buffer widths sufficient to
guarantee adequate residence times across hydrogeologic settings
\parencite{Puckett2005}, low-rate effluent irrigation to reduce direct transfers
through mole-pipe drainage \parencite{Monaghan2007}, and long-term monitoring of
tillage and earthworm effects on phosphorus leaching \parencite{Ulén2018}.
Collectively, these companion-practice frontiers reinforce the section's central
message: moving edge-of-field treatment from site-scale demonstration toward
validated, transferable design guidance capable of meeting watershed-scale goals
across diverse climates, soils, and hydrologic regimes.